% Generated draft: requires mathematical and source review.
\documentclass[11pt]{article}
\usepackage[margin=25mm]{geometry}
\usepackage{fontspec}
\setmainfont{Latin Modern Roman}
\usepackage{amsmath,amssymb,mathtools,unicode-math}
\setmathfont{Latin Modern Math}
\usepackage{enumitem,needspace,xurl,hyperref}
\setlength{\parindent}{0pt}
\setlength{\parskip}{5pt}
\setlength{\emergencystretch}{3em}
\newcommand{\N}{\mathbb{N}}
\newcommand{\Z}{\mathbb{Z}}
\newcommand{\Q}{\mathbb{Q}}
\newcommand{\R}{\mathbb{R}}
\newcommand{\C}{\mathbb{C}}
\newcommand{\F}{\mathbb{F}}
\newcommand{\A}{\mathbb{A}}
\newcommand{\PP}{\mathbb{P}}
\newcommand{\E}{\mathbb{E}}
\newcommand{\eps}{\varepsilon}
\newcommand{\dd}{\,\mathrm{d}}
\begin{document}
\section*{UnsolvedMath ID: 30003387}
\textbf{Source problem code:} catalogue entry
\subsection{Definitions and mathematical statement}
Let \(d\) be an integer with \(d>2\). The supplied record concerns the degenerate-diffusion Keller--Segel equation
\[
\partial_t\rho=\Delta(\rho^m)-\nabla\cdot\bigl(\rho\nabla(N*\rho)\bigr).
\]
Here \(t\) denotes time, \(\nabla\), \(\nabla\cdot\), and \(\Delta\) are spatial differential operators, \(m\) is the diffusion exponent, and \(N*\rho\) denotes spatial convolution with an interaction kernel \(N\). The record does not specify the spatial domain, the admissible range of \(m\), or the definition, normalization, or sign convention of \(N\). These are unresolved parameters of the formulation; in particular, no particular Newtonian kernel or exponent range is imposed here.

\textbf{Requested mathematical determination.} Determine the asymptotic behavior as \(t\to\infty\) of solutions of this equation in dimension \(d>2\), including whether they converge to stationary states or whether mass escapes to infinity. Formally, a stationary state \(\rho_\infty\) is a time-independent solution satisfying
\[
0=\Delta(\rho_\infty^m)-\nabla\cdot\bigl(\rho_\infty\nabla(N*\rho_\infty)\bigr),
\]
with the same domain, boundary conditions, and solution interpretation as the evolution equation, once these have been specified.

The source supplies neither an initial-data class nor a solution concept, global-existence hypotheses, boundary conditions, or a topology of convergence. It also does not specify whether convergence is intended for every solution in some class, along subsequences, or modulo spatial translations. Consequently, the requested determination is an incompletely specified research problem rather than a fully quantified convergence assertion.

The phrase ``lose mass at infinity'' requires an unbounded spatial setting and a specified notion of escape. For example, on \(\R^d\), if \(\rho\geq0\) is integrable, the total mass would be \(M(t)=\int_{\R^d}\rho(t,x)\,\mathrm dx\), whereas the mass in a fixed ball would be \(\int_{B_R}\rho(t,x)\,\mathrm dx\), where \(B_R=\{x\in\R^d:|x|<R\}\). Escape can concern the distribution of mass in bounded regions even when total mass is conserved. These conditional definitions clarify the terminology only: the source does not choose a domain, specify partial versus complete escape, assert mass conservation, or state that convergence and escape exhaust all possible behaviors.
\subsection{Short English statement}
For the supplied degenerate-diffusion Keller--Segel equation in dimension \(d>2\), determine the long-time behavior, including convergence to stationary states and possible escape of mass to infinity. The exponent range, interaction kernel, domain, initial data, solution class, and meanings of convergence and mass escape are missing, so the record does not determine a unique precise asymptotic claim.
\subsection{Sources}
\begin{itemize}
\item Supplied local UnsolvedMath registry record, permanent ID \(30003387\), entitled ``Long-Time Behavior of Degenerate Keller--Segel Equations,'' in the collection ``Oberwolfach Reports Open Problems.'' Exact source problem code: \texttt{catalogue entry}. The record gives proposed year \(2017\) and supplies no individual proposer.
\item Supplied report citation: \emph{Applications of Optimal Transportation in the Natural Sciences} (2018), Oberwolfach Reports, pp.~339--416, DOI \texttt{10.4171/owr/2017/7}: \url{https://doi.org/10.4171/owr/2017/7}. The supplied background describes Oberwolfach Reports as documenting mathematical workshops and their open-problem discussions. The citation year \(2018\), proposed year \(2017\), and DOI containing \(2017\) are retained as supplied; their relationship is not resolved here.
\item Supplied literature source labeled ``Oberwolfach report'': \url{https://ems.press/content/serial-article-files/46673}.
\item Supplied literature source labeled ``Carrillo--Hittmeir--Volzone--Yao'': \url{https://doi.org/10.1007/s00222-019-00898-x}.
\item Supplied literature source labeled ``Cozzi--Gie--Kelliher'': \url{https://doi.org/10.1016/j.jmaa.2017.04.009}.
\item The supplied literature assessment was checked on \textbf{2026-08-21}. This is supplied dated curation derived from the OWR record, not an independently verified assessment. It labels the problem ``partially solved'' and reports existence of stationary minimizers and several radial or low-dimensional long-time results, while describing the full higher-dimensional asymptotic behavior as unresolved. The hypotheses and precise scope of those reported results are not supplied; time-sensitive status and cited claims require independent verification.
\end{itemize}
\end{document}
